\hypertarget{c11-smart-pointers-memory-management}{%
\chapter{C++11 SMART POINTERS \& MEMORY
MANAGEMENT}\label{c11-smart-pointers-memory-management}}

C++11 revolutionized memory management by introducing smart pointers,
which strictly define ownership semantics and automate memory
reclamation, effectively making \passthrough{\lstinline!new!} and
\passthrough{\lstinline!delete!} unnecessary in user code.

\begin{center}\rule{0.5\linewidth}{0.5pt}\end{center}

\hypertarget{the-problem-with-raw-pointers}{%
\section{1. THE PROBLEM WITH RAW
POINTERS}\label{the-problem-with-raw-pointers}}

In C++98, dynamic memory required manual management: 1. \textbf{Memory
Leaks}: Forgetting \passthrough{\lstinline!delete!}. 2. \textbf{Dangling
Pointers}: Accessing deleted memory. 3. \textbf{Double Free}: Deleting
the same memory twice. 4. \textbf{Exception Safety}: If an exception
throws before \passthrough{\lstinline!delete!}, memory leaks.

Smart pointers solve these by using \textbf{RAII (Resource Acquisition
Is Initialization)}.

\begin{center}\rule{0.5\linewidth}{0.5pt}\end{center}

\hypertarget{unique_ptr-exclusive-ownership}{%
\section{2. UNIQUE\_PTR (Exclusive
Ownership)}\label{unique_ptr-exclusive-ownership}}

\passthrough{\lstinline!std::unique\_ptr!} represents exclusive
ownership. An object can have only one
\passthrough{\lstinline!unique\_ptr!} pointing to it. When the
\passthrough{\lstinline!unique\_ptr!} is destroyed, the object is
deleted.

\hypertarget{basic-usage}{%
\subsection{2.1 Basic Usage}\label{basic-usage}}

\begin{lstlisting}[language={C++}]
#include <memory>

void func() {
    std::unique_ptr<int> ptr(new int(10));
    // or better: auto ptr = std::make_unique<int>(10); (C++14)
    
    *ptr = 20;
    // No delete needed. Memory freed when ptr goes out of scope.
}
\end{lstlisting}

\hypertarget{move-only}{%
\subsection{2.2 Move Only}\label{move-only}}

You cannot copy a \passthrough{\lstinline!unique\_ptr!}. You must
\textbf{move} it. This ensures uniqueness.

\begin{lstlisting}[language={C++}]
std::unique_ptr<int> p1(new int(5));
// std::unique_ptr<int> p2 = p1; // Error! Copy deleted.

std::unique_ptr<int> p2 = std::move(p1); // OK. p1 is now empty/null.
\end{lstlisting}

\hypertarget{custom-deleters}{%
\subsection{2.3 Custom Deleters}\label{custom-deleters}}

Useful for managing C-style resources (files, sockets).

\begin{lstlisting}[language={C++}]
auto deleter = [](FILE* f) { fclose(f); };
std::unique_ptr<FILE, decltype(deleter)> file(fopen("test.txt", "r"), deleter);
\end{lstlisting}

\begin{center}\rule{0.5\linewidth}{0.5pt}\end{center}

\hypertarget{shared_ptr-shared-ownership}{%
\section{3. SHARED\_PTR (Shared
Ownership)}\label{shared_ptr-shared-ownership}}

\passthrough{\lstinline!std::shared\_ptr!} allows multiple pointers to
own the same resource. The resource is deleted only when the \emph{last}
\passthrough{\lstinline!shared\_ptr!} is destroyed.

\hypertarget{reference-counting}{%
\subsection{3.1 Reference Counting}\label{reference-counting}}

It maintains a ``control block'' with a reference count.

\begin{lstlisting}[language={C++}]
auto p1 = std::make_shared<int>(100); // Ref count = 1
{
    auto p2 = p1; // Copy allowed. Ref count = 2
} // p2 destroyed. Ref count = 1

// p1 destroyed. Ref count = 0. Memory freed.
\end{lstlisting}

\hypertarget{performance-cost}{%
\subsection{3.2 Performance Cost}\label{performance-cost}}

\passthrough{\lstinline!shared\_ptr!} is heavier than
\passthrough{\lstinline!unique\_ptr!} (2x size usually, plus atomic
ref-count increment/decrement overhead). Use only when ownership is
truly shared.

\begin{center}\rule{0.5\linewidth}{0.5pt}\end{center}

\hypertarget{weak_ptr-non-owning-reference}{%
\section{4. WEAK\_PTR (Non-Owning
Reference)}\label{weak_ptr-non-owning-reference}}

\passthrough{\lstinline!std::weak\_ptr!} observes a
\passthrough{\lstinline!shared\_ptr!} without keeping it alive. It
breaks \textbf{circular references}.

\hypertarget{circular-reference-problem}{%
\subsection{4.1 Circular Reference
Problem}\label{circular-reference-problem}}

If A has a \passthrough{\lstinline!shared\_ptr!} to B, and B has a
\passthrough{\lstinline!shared\_ptr!} to A, the reference count never
drops to zero.

\hypertarget{using-weak_ptr}{%
\subsection{4.2 Using weak\_ptr}\label{using-weak_ptr}}

\begin{lstlisting}[language={C++}]
struct B;
struct A {
    std::shared_ptr<B> b_ptr;
};
struct B {
    std::weak_ptr<A> a_ptr; // Use weak_ptr back to A
};
\end{lstlisting}

To use a \passthrough{\lstinline!weak\_ptr!}, you must convert it to
\passthrough{\lstinline!shared\_ptr!} via
\passthrough{\lstinline!.lock()!}.

\begin{lstlisting}[language={C++}]
if (auto shared = weak.lock()) {
    // safe to use shared
} else {
    // object died
}
\end{lstlisting}

\begin{center}\rule{0.5\linewidth}{0.5pt}\end{center}

\hypertarget{best-practices}{%
\section{5. BEST PRACTICES}\label{best-practices}}

\begin{enumerate}
\def\labelenumi{\arabic{enumi}.}
\tightlist
\item
  \textbf{Prefer \passthrough{\lstinline!unique\_ptr!}} by default. It
  has zero overhead.
\item
  \textbf{Use \passthrough{\lstinline!make\_unique!}} (C++14) and
  \textbf{\passthrough{\lstinline!make\_shared!}}. They are cleaner and
  exception-safe. \passthrough{\lstinline!make\_shared!} is also more
  efficient (allocates object and control block in one chunk).
\item
  \textbf{Avoid \passthrough{\lstinline!new!} and
  \passthrough{\lstinline!delete!}}.
\end{enumerate}
